\chapter{Results}
\label{chap:results}

This chapter presents the results by subsystem, followed by a full-system test on the vehicle, and maps each back to the research questions and requirements in Table~\ref{tab:requirements-mapping}. Bench-level validation results are reported alongside their methods in Chapter~\ref{chap:method}.

\section{Inertial acquisition integrity (RQ3)}
Achieved sample rate, inter-sample jitter distribution, and dropped-sample counts for both IMUs, first on the bench and then on the moving vehicle.

\section{Sensor calibration}
Per-sensor magnitude calibration factors and the resulting agreement between the two IMUs, at rest and under dynamic excitation.

\section{Ride characterisation (RQ1, RQ4)}
Time-domain features (RMS, peak-to-peak, crest factor, jerk, threshold crossings) and the speed-normalised roughness index recorded over known road sections, compared against identifiable road and powertrain events.

\section{Battery-management telemetry}
BLE connection reliability over a sustained session, and a measured discharge curve compared against the pack's rated capacity.

\section{Camera, synchronisation and proximity sensing}
Camera-to-IMU time alignment, object-detection alert events and the fields recorded for them, and time-of-flight distance measurements, including behaviour against light and dark targets.

\section{Cellular synchronisation and data integrity (RQ3)}
Publish success rate, latency, and recovery behaviour after induced or naturally occurring loss of coverage, with local-log completeness as the reference.

\section{Driver display (RQ2)}
Information hierarchy of the dashboard page as implemented, alert update latency to the screen, legibility under representative lighting, and observed glance time.

\section{Full-system test}
All subsystems operating together on the vehicle over a defined route.
